\usepackage{listings}

\DeclareMathOperator{\rightleftarrow}{\rightarrow\!\leftarrow}
\let\implies\relax
\DeclareMathOperator{\implies}{\rightarrow}
\let\iff\relax
\DeclareMathOperator{\iff}{\leftrightarrow}
\newcommand{\fn}[2]{#1\hspace{-0.1em}\left(#2\right)}
\newcommand{\sfrac}[2]{{^{#1\hspace{-0.1em}}/_{\hspace{-0.1em}#2}}}
\newcommand{\bset}[2]{\left\{ #1 \:\left\lvert\; #2 \right.\right\}}
\newcommand{\dplus}{\mathbin{+\hspace{-0.6em}+}}
\newcommand{\np}{\mathcal{N\!P}}
\newcommand{\p}{\mathcal{P}}

% TODO: This is known to be a temporary workaround for tagging listings that
%       may not work in later releases or in all possible use cases. See
%       https://github.com/latex3/tagging-project/issues/70 and
%       https://tex.stackexchange.com/questions/755794/.
\makeatletter
\ExplSyntaxOn
\DeclareInstance{blockenv}{lstlisting}{display}{
    name = lstlisting,
    tag-name = verbatim,
    tag-attr-class = ,
    tagging-recipe = standard,
    inner-level-counter = ,
    increment-level = false,
    setup-code = ,
    block-instance = displayblock ,
    inner-instance = ,
    final-code = \tl_set:Nn \l__tag_para_main_tag_tl {codeline}\tagtool{paratag=Code},
}
\ExplSyntaxOff
\lstnewenvironment{taglisting}[2][]{%
    \UseInstance{blockenv}{lstlisting} {}
    \lst@TestEOLChar{#2}%
    \lstset{#1}%
    \csname\@lst @SetFirstNumber\endcsname%
}{%
    \@nobreakfalse
    \csname\@lst @SaveFirstNumber\endcsname%
    \endblockenv
}
\makeatother

\lstset{
    numbers=left,
    numberstyle=\scriptsize\ttfamily,
    basicstyle=\normalsize\ttfamily,
    tabsize=4,
    frame=L,
    xleftmargin=1.75em,
    aboveskip=0.4375em,
    belowskip=0em,
    upquote=true,
    showlines=true,
    basewidth=0.5625em
}

